\subsection{EQUIVALENT MATRICES; SIMILAR MATRICES}

DEFINITION 5. Two matrices \(X, X'\) with \(m\) rows and \(n\) columns over a ring are called equivalent if there exists an invertible square matrix \(P\) of order \(m\) and an invertible square matrix \(Q\) of order \(n\) such that

\[
X ^ {\prime} = P X Q.\tag{40}
\]

Clearly the relation ``X and \(X'\) are equivalent'' is an equivalence relation (Set Theory, II, § 6, no. 1) on the set \(A^{mn}\) of matrices of type \((m, n)\) over A, which justifies the terminology.

With this definition, Proposition 6 of no. 8 can be stated by saying that when the bases are changed in two right A-modules E, F (with finite bases), the matrix of a linear mapping \(u \colon \mathbf{E} \to \mathbf{F}\) with respect to the new bases is equivalent to the matrix of \(u\) with respect to the old bases.

Conversely, if relation (40) holds and \(u \colon \mathbf{A}_d^n \to \mathbf{A}_d^m\) is a linear mapping whose matrix is \(X\) with respect to the respective canonical bases \((e_i)\) and \((f_j)\) of \(\mathbf{A}_d^n\) and \(\mathbf{A}_d^m\) , then \(X'\) is the matrix of \(u\) with respect to the bases \((e_i')\) and \((f_j')\) such that \(Q\) is the matrix of passage from \((e_i)\) to \((e_i')\) and \(P^{-1}\) the matrix of passage from \((f_j)\) to \((f_j')\) .

Examples of equivalent matrices. (1) Two matrices \(X = (x_{ij})\) and \(X' = (x_{ij}')\) with \(m\) rows and \(n\) columns ``differ only in the order of their rows'' if there exists a permutation \(\sigma\) of the interval \([1, m]\) of \(\mathbf{N}\) , such that for every ordered pair of indices \((i,j)\) , \(x_{ij}' = x_{\sigma(i),j}\) (we also say that \(X'\) is obtained by performing the permutation \(\sigma^{-1}\) on the rows of \(X\) ). The matrices \(X\) and \(X'\) are then equivalent, for \(X' = PX\) , where \(P\) is the matrix of the permutation \(\sigma^{-1}\) (cf.~no. 7, Example II).

Similarly X and \(X'\) are said to differ only in the order of their columns if there exists a permutation \(\tau\) of \([1, n]\) such that \(x_{ij}' = x_{i,\tau(j)}\) for every ordered pair of indices \((i,j)\) , X and \(X'\) are also equivalent, for \(X' = XQ\) where Q is the matrix of the permutation \(\tau\) .

Note that in the above notation \(P\) is the matrix of passage from a basis \((f_j)_{1 \leqslant j \leqslant m}\) to the basis \((f_\sigma^{-1}(j))_{1 \leqslant j \leqslant m}\) and \(Q\) the matrix of passage from a basis \((e_i)_{1 \leqslant i \leqslant n}\) to the basis \((e_{\tau(i)})_{1 \leqslant i \leqslant n}\) .

\begin{enumerate}
\def\labelenumi{(\arabic{enumi})}
\setcounter{enumi}{1}
\tightlist
\item
  Let \(j, k\) be distinct elements of \([1, n]\) and let \(a \in A\) .
\end{enumerate}

Suppose that for \(1 \leqslant i \leqslant m\) , \(x_{ij}' = x_{ij} + x_{ik}a\) and \(x_{il}' = x_{il}\) for \(j \neq l\) and \(1 \leqslant i \leqslant m\) ; \(X'\) is said to be derived from \(X\) by adding to the column of \(X\) of index \(j\) the column of index \(k\) multiplied on the right by \(a\) . In this case \(X\) and \(X'\) are also equivalent: for if \(Q = I_n + aE_{kj}\) (an invertible triangular matrix, as seen in no. 7), then \(X' = XQ\) .

Similarly, let \(h, i\) be two distinct elements of \([1, m]\) and \(a\) an element of \(A\) ; if \(X'\) is derived from \(X\) by adding to the row of \(X\) of index \(i\) the row of index \(h\) multiplied on the left by \(a, X\) and \(X'\) are equivalent, for \(X' = PX\) , where \(P = I_m + aE_{th}\) .

\begin{enumerate}
\def\labelenumi{(\arabic{enumi})}
\setcounter{enumi}{2}
\tightlist
\item
  Finally, if, for a given index \(j\) , \(x_{ij}' = x_{ij}c\) for \(1 \leqslant i \leqslant m\) , where \(c\) is invertible and \(x_{il}' = x_{il}\) for \(1 \leqslant i \leqslant m\) and \(l \neq j\) , \(X\) and \(X'\) are equivalent; for \(X' = XQ\) , where \(Q\) is the matrix \(\text{diag}(a_k)\) with \(a_j = c\) , \(a_k = 1\) for \(k \neq j\) . Then \(X'\) is said to be derived from \(X\) by multiplying the column of \(X\) of index \(j\) on the right by \(a\) .
\end{enumerate}

Similarly, if \(X'\) is derived from \(X'\) by multiplying the row of \(X\) of index \(i\) on the left by an invertible element \(c \in A\) , \(X'\) and \(X\) are equivalent, for \(X' = PX\) where \(P\) is the matrix \(\text{diag}(b_h)\) with \(b_i = c\) , \(b_h = 1\) for \(h \neq i\) .

DEFINITION 6. Two square matrices \(X, X'\) of order \(n\) over a ring \(A\) are called similar if there exists an invertible square matrix \(P\) of order \(n\) such that

\[
X ^ {\prime} = P X P ^ {- 1}\tag{41}
\]

Clearly the relation ``X and \(X'\) are similar'' is an equivalence relation on \(\mathbf{M}_{n}(\mathbf{A})\) meaning that X and \(X'\) are transformed into one another by an inner automorphism of this ring.

With this definition, Corollary 1 to Proposition 6 of no. 8 can be stated by saying that when the basis of an A-module E (with a finite basis) is changed, the matrix of an endomorphism u with respect to the new basis is similar to the matrix of u with respect to the old basis.

Remarks. (1) Two square matrices which differ only in the order of their rows (or the order of their columns) are equivalent, but in general not similar. A matrix similar to a square matrix \(X = (x_{ij})\) can be obtained by performing the same permutation \(\sigma^{-1}\) on the rows and columns, that is by considering the matrix \(X' = (x_{ij}')\) , where \(x_{ij}' = x_{\sigma(i),\sigma(j)}\) for every ordered pair of indices; for if \(X\) is the matrix of an endomorphism \(u\) of \(A_d^n\) with respect to a basis \((e_i)_{1 \leq i \leq n}\) , \(X'\) is the matrix of \(u\) with respect to the basis \((e_{\sigma(i)})_{1 \leq i \leq n}\) .

\begin{enumerate}
\def\labelenumi{(\arabic{enumi})}
\setcounter{enumi}{1}
\tightlist
\item
  Let \(X\) and \(X'\) be two square matrices of order \(n\) which can be written in the form of diagonal matrices of square matrices (no. 7, Example IV):
\end{enumerate}

\[
X = \left( \begin{array}{c c c c} X _ {1} & 0 & \ldots & 0 \\ 0 & X _ {2} & \ldots & 0 \\ \cdot & \cdot & \cdot & \cdot \\ 0 & 0 & \ldots & X _ {p} \end{array} \right) \qquad X ^ {\prime} = \left( \begin{array}{c c c c} X _ {1} ^ {\prime} & 0 & \ldots & 0 \\ 0 & X _ {2} ^ {\prime} & \ldots & 0 \\ \cdot & \cdot & \cdot & \cdot \\ 0 & 0 & \ldots & X _ {p} ^ {\prime} \end{array} \right)
\]

corresponding to the same partition of the indexing set \([1, n]\) for X and \(X'\) . If, for \(1 \leqslant i \leqslant p\) , \(X_i\) and \(X'_i\) are equivalent (resp. similar), then X and \(X'\) are equivalent (resp. similar): for, if \(X'_i = P_i X_i Q_i\) for \(1 \leqslant i \leqslant p\) , then \(X' = PXQ\) where

\[
P = \left( \begin{array}{c c c c c} P _ {1} & 0 & \dots & 0 \\ 0 & P _ {2} & \dots & 0 \\ \cdot & \cdot & \cdot & \cdot & \cdot \\ 0 & 0 & \dots & P _ {p} \end{array} \right) \quad Q = \left( \begin{array}{c c c c c} Q _ {1} & 0 & \dots & 0 \\ 0 & Q _ {2} & \dots & 0 \\ \cdot & \cdot & \cdot & \cdot & \cdot \\ 0 & 0 & \dots & Q _ {p} \end{array} \right)
\]

as follows from calculating the ``block'' product (no. 6). Moreover, if \(Q_{i} = P_{i}^{-1}\) for all i, then \(Q = P^{-1}\) .

